\documentclass[11pt,a4paper]{article}
\usepackage[utf8]{inputenc}
\usepackage[spanish]{babel}
\usepackage{hyperref}
\usepackage{geometry}
\geometry{top=2.5cm, bottom=2.5cm, left=2.5cm, right=2.5cm}
\usepackage{graphicx} 
\usepackage{tikz}
\usetikzlibrary{shapes.geometric, arrows.meta, positioning, fit, backgrounds, babel}
\usepackage{float}
\usepackage{xcolor}

% Definición de paleta de colores para el diagrama
\definecolor{colorSource}{HTML}{D5E8D4} % Verde suave
\definecolor{colorBeam}{HTML}{FFE6CC}   % Naranja suave
\definecolor{colorError}{HTML}{F8CECC}  % Rojo suave
\definecolor{colorSilver}{HTML}{E1D5E7} % Morado suave
\definecolor{colorGold}{HTML}{FFF2CC}   % Amarillo suave
\definecolor{colorAirflow}{HTML}{DAE8FC} % Azul suave

\title{Proyecto Final: Sistema de Consolidación de Ventas Omnicanal Super-Market}
\author{Geraldine Cornejo, Javiera Barrales, Rodrigo Flores}
\date{Junio 2026}

\begin{document}

\maketitle

\section{Introducción al Proyecto}

La corporación multinacional Super-Market enfrenta un desafío crítico de fragmentación de datos: las operaciones de Santiago, Lima y Buenos Aires generan información local en archivos CSV que presentan inconsistencias de formato, anomalías financieras y disparidad en las monedas. 

Para proveer un repositorio de confianza al equipo de Data Science, implementamos una arquitectura Medallion. Esta solución centraliza y denormaliza las transacciones mediante Apache Beam (Capa Silver en Parquet), para luego modelar y estandarizar métricas financieras a USD mediante dbt en PostgreSQL (Capa Gold), bajo la orquestación resiliente de Apache Airflow.

\section{Blueprint de Arquitectura y Flujo de Datos}

A continuación, se presenta el Diagrama de Flujo de Datos (DFD) que ilustra el ciclo de vida de la información de forma segmentada.

\begin{figure}[H]
\centering
\begin{tikzpicture}[
    node distance=1.2cm and 1cm,
    base/.style={rectangle, rounded corners, draw=black, thick, align=center, font=\scriptsize, inner sep=1ex},
    source/.style={base, fill=colorSource, text width=2.4cm},
    beam/.style={base, fill=colorBeam, text width=2.6cm},
    error/.style={base, fill=colorError, text width=2.4cm},
    storage/.style={cylinder, shape border rotate=90, draw=black, thick, aspect=0.25, align=center, font=\scriptsize, inner sep=1ex},
    silver/.style={storage, fill=colorSilver, text width=2.6cm},
    gold/.style={storage, fill=colorGold, text width=2.6cm},
    airflow/.style={base, fill=colorAirflow, text width=12cm, font=\bfseries\small},
    flecha/.style={thick, ->, >=stealth},
    flecha_texto/.style={font=\tiny, align=center}
]

% Orquestación
\node (airflow) [airflow] {Apache Airflow (DAG, Dependencias y Retries)};

% Capa de Origen (Sources)
\node (sales) [source, below left=1.2cm and -2cm of airflow] {sales\_data.csv\\(Ventas)};
\node (logs) [source, below right=1.2cm and -2cm of airflow] {status\_logs.csv\\(Historial)};

% Procesamiento (Apache Beam)
\node (beam_read) [beam, below=1.8cm of airflow] {Lectura Asíncrona\\(Multithreading)};
\node (join) [beam, below=1cm of beam_read] {Join: CoGroupByKey\\(Por transaction\_id)};
\node (rejected) [error, right=1.5cm of join] {Side Output\\(rejected\_sales.csv)};

% Capa de Persistencia (Silver)
\node (parquet) [silver, below=1.5cm of join] {Data Lake\\(Parquet Anidado)};

% Capa de Transformación (dbt) y Servicio (Gold)
\node (dbt) [beam, below=1.5cm of parquet] {dbt run \& test\\(Filtros y Conversión)};
\node (postgres) [gold, below=1.5cm of dbt] {PostgreSQL\\(fct\_sales\_daily y KPIs)};

% Conexiones Sources -> Beam
\draw [flecha] (sales) |- node[flecha_texto, above, pos=0.75] {Validación} (beam_read);
\draw [flecha] (logs) |- node[flecha_texto, above, pos=0.75] {Validación} (beam_read);

% Conexiones Beam interno
\draw [flecha] (beam_read) -- (join);
\draw [flecha] (join) -- node[flecha_texto, above] {Anomalías} (rejected);

% Conexiones Beam -> Silver
\draw [flecha] (join) -- node[flecha_texto, right] {PyArrow Schema} (parquet);

% Conexiones Silver -> dbt -> Gold
\draw [flecha] (parquet) -- (dbt);
\draw [flecha] (dbt) -- (postgres);

% Conexiones Airflow de control (Líneas punteadas)
\draw [dashed, ->, thick, color=blue!70] (airflow.south west) -- ++(0,-0.8) -| node[flecha_texto, left, pos=0.9] {Dispara Task 1} (beam_read.north west);
\draw [dashed, ->, thick, color=blue!70] (airflow.south east) -- ++(0,-0.8) -| node[flecha_texto, right, pos=0.9] {Dispara Task 3} (dbt.north east);

% Fondos agrupadores con etiquetas sutiles (\tiny y gray!80)
\begin{scope}[on background layer]
    \node [draw=gray, dashed, thick, fill=gray!5, fit=(sales)(logs), inner sep=0.4cm, label={[font=\tiny, text=gray!80]above left:Capa de Origen (Sources)}] {};
    \node [draw=gray, dashed, thick, fill=gray!5, fit=(beam_read)(join)(rejected), inner sep=0.4cm, label={[font=\tiny, text=gray!80]above left:Procesamiento (Apache Beam)}] {};
    \node [draw=gray, dashed, thick, fill=gray!5, fit=(parquet), inner sep=0.4cm, label={[font=\tiny, text=gray!80]above left:Capa de Persistencia (Silver)}] {};
    \node [draw=gray, dashed, thick, fill=gray!5, fit=(dbt), inner sep=0.4cm, label={[font=\tiny, text=gray!80]above left:Capa de Transformación (dbt)}] {};
    \node [draw=gray, dashed, thick, fill=gray!5, fit=(postgres), inner sep=0.4cm, label={[font=\tiny, text=gray!80]above left:Capa de Servicio (Gold Serving)}] {};
\end{scope}

\end{tikzpicture}
\caption{Diagrama de Arquitectura segmentado por Capas Técnicas}
\label{fig:dfd}
\end{figure}

\clearpage

\section{Análisis Técnico y Decisiones Arquitectónicas}

\subsection{Contratos de Esquema y Calidad}
Para proteger la integridad del dato desde su origen débil (CSV) hasta la capa analítica, definimos contratos estrictos. En Apache Beam, la estructura denormalizada se serializa utilizando el objeto \texttt{SILVER\_SCHEMA} de PyArrow. Al pasarlo a la función \texttt{WriteToParquet}, PyArrow actúa como un guardián (gatekeeper) que rechaza cualquier escritura que no cumpla con los tipos y jerarquías anidadas. 

En la Capa Gold, la calidad se garantiza matemáticamente mediante \texttt{schema.yml} en dbt. Aplicamos pruebas genéricas (\texttt{unique} y \texttt{not\_null} sobre \texttt{transaction\_id}) y usamos el paquete \texttt{dbt\_expectations} para forzar que el \texttt{amount\_usd} tenga un límite rígido de \texttt{min\_value: 0.01}, asegurando que ninguna tasa fraccionaria anómala afecte los KPIs del negocio.

\subsection{Orquestación del DAG}
El grafo de Airflow se configuró secuencialmente: \texttt{Beam $\rightarrow$ Sensor $\rightarrow$ dbt}. Apoyándonos en la regla \texttt{all\_success}, garantizamos que dbt se detenga por completo si los pasos anteriores fallan. 
Para aumentar la resiliencia, configuramos la tarea con 2 reintentos (\texttt{retries}) y un \texttt{retry\_delay} de 5 minutos. Además, implementamos un \texttt{FileSensor} que hace *poking* cada 30 segundos verificando la creación exitosa del Parquet en la Capa Silver. Si transcurren 600 segundos (10 minutos), el sensor lanza un timeout y aborta el pipeline, evitando que dbt procese datos corruptos o incompletos.

\subsection{Matriz de Resiliencia y Procesamiento Paralelo}
Para la ingesta inicial, aprovechamos las capacidades distribuidas nativas de Apache Beam (específicamente mediante \texttt{beam.io.ReadFromText}), lo que nos permite un procesamiento óptimo y en paralelo de los chunks de texto sin la necesidad de implementar bloqueos manuales. Esto garantiza una lectura eficiente y libre de contención.

Sobre esta lectura paralela, aplicamos las reglas de validación: las ventas exigen 5 columnas y un \texttt{amount} estricto $>0$ (casteado a float); los logs exigen 3 columnas y un dominio cerrado (\texttt{CREATED, PENDING, COMPLETED, REFUNDED}). Si una regla falla, el registro emite un \texttt{beam.pvalue.TaggedOutput('rejected', ...)}. Posteriormente, unificamos estos flujos de error con \texttt{beam.Flatten()} y los enviamos al archivo \texttt{rejected\_sales.csv} para auditoría, manteniendo intacto el pipeline principal.

\subsection{Sinking Strategy e Idempotencia}
El uso de Apache Parquet permite compresión masiva y aislamiento por columnas. Para asegurar la idempotencia del pipeline en Beam, configuramos la escritura \texttt{WriteToParquet} sin el flag de *append*, forzando una sobreescritura limpia en cada ejecución diaria y estampando la metadata \texttt{processed\_at}. 
En paralelo, la idempotencia en dbt (Capa Gold) se garantizó definiendo explícitamente \texttt{materialized='table'} en \texttt{fct\_sales.sql}, lo que instruye a PostgreSQL a eliminar y recrear la tabla completa en cada corrida, eliminando el riesgo de duplicidad por *upserts* fallidos.

\subsection{El Modelo de las 4 Preguntas de Beam}
\begin{itemize}
    \item \textbf{What (Qué se calcula):} Join de ventas e histórico de estados, denormalizándolos mediante tipado PyArrow estricto, aislando los datos anómalos.
    \item \textbf{Where (Dónde en el tiempo):} Uso de Ventanas Globales para el procesamiento Batch sobre datos limitados.
    \item \textbf{When (Cuándo se procesa):} Al finalizar el cierre contable diario.
    \item \textbf{How (Cómo se relacionan):} Usamos \texttt{CoGroupByKey} para agrupar las streams por \texttt{transaction\_id}. Luego, la función \texttt{BuildNestedRecordFn} itera sobre estos agrupamientos, parseando las fechas string a objetos \texttt{datetime} UTC y anexándolos como diccionarios dentro de la lista \texttt{status\_history}.
\end{itemize}

\subsection{Repositorio}
Nuestra solución se encuentra en:
\href{https://github.com/trencitos/PROYECTO_BD_A_GRAN_ESCALA.git}{Ir a GitHub}

\end{document}